%!TEX root = ../main_thesis.tex

\chapter{Analisi semantica: tipi e trasformazione del modello}
\label{cap:semantica}

Una volta che il sorgente è stato tradotto nella rappresentazione intermedia
(Capitolo~\ref{cap:parsing}), il compilatore deve verificarne la correttezza e
trasformarlo in una forma concreta. Queste due attività costituiscono l'analisi
semantica.

\section{Concetti teorici}
\label{sec:semantica-teoria}

\paragraph{Analisi semantica.}
Mentre l'analisi sintattica stabilisce se un testo è \emph{ben formato},
l'analisi semantica stabilisce se esso è anche \emph{dotato di significato}
secondo le regole del linguaggio. Sono errori semantici, e non sintattici, l'uso
di una variabile non dichiarata, l'applicazione di un operatore a operandi
incompatibili o l'invocazione di una funzione con un numero errato di argomenti.

\paragraph{Sistemi di tipi.}
Lo strumento principale con cui si esprimono e si verificano queste regole è il
\emph{sistema dei tipi}~\cite{pierce2002}. Un tipo classifica i valori in base
alla loro natura (numeri, stringhe, collezioni, e così via) e impone vincoli su
come essi possono essere combinati. La verifica statica dei tipi (\emph{type
checking}) controlla tali vincoli prima dell'esecuzione, permettendo di
respingere i programmi scorretti con un messaggio diagnostico anziché lasciarli
fallire in un momento successivo. È utile distinguere fra il \emph{valore} che
un'espressione assume e il \emph{tipo} che la descrive: il type checker ragiona
sui tipi, ignorando i valori concreti.

\paragraph{Ambiti e tabella dei simboli.}
Per sapere a cosa si riferisce un nome, il type checker mantiene una
\emph{tabella dei simboli} che associa a ogni identificatore il suo tipo. Poiché
i linguaggi prevedono costrutti che introducono nomi validi solo localmente
(in ROOC, gli iteratori), la tabella è organizzata in \emph{ambiti} (\emph{scope})
annidati, tipicamente gestiti come una pila: l'ingresso in un nuovo costrutto
introduce un ambito, la sua uscita lo rimuove, e la ricerca di un nome procede
dall'ambito più interno verso quelli più esterni.

\section{Applicazione in ROOC: il sistema dei tipi}
\label{sec:tipi-rooc}

In ROOC ogni valore manipolato a tempo di esecuzione è una \texttt{Primitive}
(Capitolo~\ref{cap:linguaggio}); a essa corrisponde, sul piano dei tipi, una
\texttt{PrimitiveKind}, che ne è la controparte statica. Mentre una
\texttt{Primitive} contiene un valore concreto, una \texttt{PrimitiveKind} ne
descrive soltanto la natura (Listato~\ref{lst:primitivekind}).

\

\begin{lstlisting}[language=rust, caption={Versione semplificata dell'enumerazione \texttt{PrimitiveKind}.}, label={lst:primitivekind}]
enum PrimitiveKind {
    Number, Integer, PositiveInteger,
    String, Boolean,
    Graph, GraphNode, GraphEdge,
    Iterable(Box<PrimitiveKind>),
    Tuple(Vec<PrimitiveKind>),
    Any,
    // ...
}
\end{lstlisting}

Due varianti meritano attenzione. \texttt{Iterable} è \emph{parametrica} sul tipo
dei propri elementi: il tipo di \texttt{[1, 2, 3]} non è semplicemente ``collezione'',
ma ``collezione di interi''. Analogamente, \texttt{Tuple} memorizza la sequenza dei
tipi delle proprie componenti, eventualmente eterogenei. La variante \texttt{Any}
rappresenta infine un tipo compatibile con qualunque altro, utile nei casi in cui
una verifica troppo rigida sarebbe inopportuna.

Il legame tra espressioni e tipi è dato dal \emph{trait} \texttt{WithType}, che
ogni espressione implementa: il suo metodo \texttt{get\_type} calcola la
\texttt{PrimitiveKind} di un'espressione a partire dal contesto. Su questa base
si esprimono le regole del linguaggio: l'applicabilità di un operatore a una
coppia di tipi e la corrispondenza tra gli argomenti di una chiamata e la
\emph{firma} della funzione invocata. La violazione di tali regole produce un
errore (per esempio un errore di operatore quando i due operandi hanno tipi
incompatibili), trattato nella Sezione~\ref{sec:trasformazione-modello}.

\section{Il type checker}
\label{sec:type-checker}

La verifica dei tipi è realizzata dal trait \texttt{TypeCheckable}, implementato
da tutte le entità dell'IR (espressioni, vincoli, dichiarazioni, l'intero
modello). Il suo metodo \texttt{type\_check} attraversa ricorsivamente la
struttura, segnalando il primo errore incontrato.

\paragraph{Il contesto di verifica.}
Lo stato della verifica è mantenuto da un \texttt{TypeCheckerContext}, che
riunisce tre informazioni: una pila di \emph{frame}, il \emph{dominio statico}
e una mappa dei token tipati. La pila di frame realizza gli ambiti annidati
discussi in precedenza: ogni \texttt{Frame} è una tabella che associa nomi a
\texttt{PrimitiveKind}, l'ingresso in un iteratore aggiunge un frame e la sua
uscita lo rimuove. La ricerca del tipo di un nome scorre i frame dal più interno
al più esterno, riproducendo le regole di visibilità lessicale. Il dominio
statico raccoglie invece i tipi delle variabili decisionali dichiarate nel
blocco \texttt{define}.

\paragraph{Dichiarazione dei nomi.}
Quando un iteratore introduce una nuova variabile, questa viene dichiarata nel
frame corrente. La dichiarazione rifiuta i nomi che coincidono con parole
riservate e, in modalità stretta, segnala la ridichiarazione di un nome già
esistente. Regole
analoghe valgono per gli indici delle variabili composte, che devono avere un
tipo ammissibile (numerico, stringa o nodo di grafo) per poter formare il nome di
una variabile del modello.

\paragraph{La mappa dei token tipati.}
Oltre a verificare la correttezza, il contesto popola una \emph{mappa dei token}
che associa a ogni posizione del sorgente il tipo dell'elemento corrispondente,
sotto forma di \texttt{TypedToken} (posizione, tipo ed eventuale identificatore).
Questa mappa non è necessaria alla compilazione in sé, ma è ciò che permette
all'editor della piattaforma web di mostrare il tipo di un'espressione al
passaggio del cursore e di fornire suggerimenti, come si vedrà nel
Capitolo~\ref{cap:piattaforma}.

\paragraph{Un esempio di errore.}
Per rendere concreto il comportamento del type checker, si supponga di alterare il
modello dello zaino (Listato~\ref{lst:knapsack}) passando alla funzione
\texttt{enumerate} la costante scalare \texttt{capacity} (un \texttt{Integer})
anziché un array. Poiché \texttt{enumerate} richiede una collezione, la verifica
fallisce e produce il messaggio del Listato~\ref{lst:typeerror}.

\

\begin{lstlisting}[language={}, caption={Errore di tipo con la relativa traccia posizionale.}, label={lst:typeerror}]
[WrongArgument] expected "Any[]", got "Integer"
	at 5:34 "capacity"
	at 5:24 "enumerate(capacity)"
	at 5:5 "sum((weight, i) in enumerate(capacity)) { weight * x_i }"
\end{lstlisting}

Il messaggio riporta il tipo atteso (\texttt{Any[]}, cioè una collezione di
elementi di tipo qualsiasi) e quello effettivamente ricevuto (\texttt{Integer}),
seguiti da una \emph{traccia} delle espressioni che racchiudono il punto
incriminato, dalla foglia (\texttt{capacity}, riga~5 colonna~34) fino all'intero
vincolo. Le posizioni provengono dalle annotazioni \texttt{Spanned} conservate
nell'IR (Capitolo~\ref{cap:parsing}) e consentono di indicare con precisione il
punto in cui l'errore ha origine.

\section{Trasformazione \texttt{PreModel} \texorpdfstring{$\rightarrow$}{->} \texttt{Model}}
\label{sec:trasformazione-modello}

Superata la verifica dei tipi, il \texttt{PreModel} viene trasformato in un
\texttt{Model}. Questa fase è un vero e proprio \emph{lowering}: le astrazioni di
alto livello del linguaggio vengono risolte ed espanse, producendo
un'espressione molto più semplice e concreta.

\paragraph{Preparazione del contesto.}
Il punto di ingresso raccoglie tutte le costanti disponibili (quelle predefinite,
quelle eventualmente fornite dall'esterno e quelle dichiarate nel blocco
\texttt{where}) e tutte le funzioni (predefinite e definite dall'utente), e
costruisce un \texttt{TransformerContext}. Quest'ultimo è l'analogo, sul piano
dei \emph{valori}, del contesto di verifica: i suoi frame associano i nomi a
valori concreti (\texttt{Primitive}) anziché a tipi, e un dominio raccoglie le
variabili decisionali via via generate.

\paragraph{Da \texttt{PreExp} a \texttt{Exp}.}
Il risultato della trasformazione delle espressioni è il tipo \texttt{Exp}
(Listato~\ref{lst:exp}), notevolmente più semplice della \texttt{PreExp} di
partenza. Tutti i costrutti di alto livello sono scomparsi: restano soltanto
numeri, variabili (identificate dal loro nome) e operazioni.

\begin{lstlisting}[language=rust, caption={Versione semplificata dell'enumerazione \texttt{Exp}.}, label={lst:exp}]
enum Exp {
    Number(f64),
    Variable(String),
    BinOp(BinOp, Box<Exp>, Box<Exp>),
    UnOp(UnOp, Box<Exp>),
    // ...
}
\end{lstlisting}

La semplificazione è la conseguenza di tre operazioni svolte durante la
trasformazione. Primo, le costanti e gli accessi agli array vengono
\emph{valutati}: un riferimento a una costante è sostituito dal suo valore.
Secondo, le \emph{block function} vengono \emph{espanse}: una sommatoria su un
iteratore diventa un albero di addizioni tra i termini concreti. Terzo, le variabili composte vengono risolte nelle
variabili concrete del modello.

\paragraph{Espansione degli iteratori e delle variabili composte.}
Il cuore dell'espansione è una procedura ricorsiva di risoluzione degli insiemi
di iterazione. Dato uno o più iteratori, essa valuta il rispettivo iterabile e,
per ogni combinazione dei valori, introduce un nuovo frame in cui lega le
variabili di iterazione ai valori correnti e invoca una funzione di
\emph{callback}. La stessa procedura serve in due contesti: replicare un vincolo
dotato di clausola \texttt{for}, generando un vincolo distinto per ogni
combinazione di indici, ed espandere il corpo di una \emph{block scoped function}
a ogni iterazione, combinando i risultati con l'operatore di aggregazione.

In entrambi i contesti, ogni variabile composta come \texttt{x\_i} viene tradotta
in una variabile del modello dal nome univoco: gli indici, ormai ridotti a valori
concreti, sono convertiti in stringa secondo regole semplici (un numero diventa la
sua rappresentazione decimale, un valore booleano diventa \texttt{T} o \texttt{F},
una stringa resta invariata, un nodo di grafo diventa il proprio nome) e
concatenati con un \emph{underscore}. Ogni variabile così ottenuta è registrata nel
dominio come una \texttt{DomainVariable}, che ne ricorda il tipo, la posizione di
origine e il numero di utilizzi.

I tre meccanismi (iterazione di un vincolo, espansione di una \emph{block scoped
function} e risoluzione delle variabili composte) agiscono di norma insieme. Si
consideri il vincolo del Listato~\ref{lst:espansione}.

\
 
\begin{lstlisting}[caption={Un vincolo che combina iterazione, sommatoria con \emph{scope} e variabili composte.}, label={lst:espansione}]
sum(j in 0..i + 1) { x_j } <= i + 1 for i in 0..3
\end{lstlisting}

La clausola \texttt{for i in 0..3} replica il vincolo per \texttt{i = 0, 1, 2}; per
ciascun valore di \texttt{i} la \emph{block scoped function} \texttt{sum} itera
\texttt{j} sull'intervallo \texttt{0..i + 1}, il cui estremo dipende dall'indice
esterno, espandendo la sommatoria in un albero di addizioni, e infine ogni
\texttt{x\_j} è risolta nella corrispondente variabile del modello. Il risultato è
la famiglia di vincoli concreti

\begin{equation*}
\begin{aligned}
i = 0:\quad & x_0 \le 1, \\
i = 1:\quad & x_0 + x_1 \le 2, \\
i = 2:\quad & x_0 + x_1 + x_2 \le 3.
\end{aligned}
\end{equation*}

\paragraph{Errori di trasformazione.}
Gli errori che possono emergere in questa fase (variabili non dichiarate,
operatori inapplicabili, firme di funzione non rispettate, indici di tipo errato,
e così via) sono rappresentati da un \texttt{TransformError}. Come anticipato nel
Capitolo~\ref{cap:parsing}, man mano che un errore si propaga verso l'esterno gli
viene associata la posizione del costrutto che lo ha generato, costruendo una
\emph{traccia} che, a partire dal sorgente, permette di indicare con precisione
dove l'errore ha avuto origine.
